\cleardoublepage
\vspace*{\fill}
\begin{abstract}
    This thesis investigates the design and implementation of a Small Language Model (SLM) using a hybrid architecture that integrates Transformer \textit{multi-head attention} mechanisms with Gated Recurrent Units (GRU). Rather than relying on purely attention-based designs, the proposed approach combines global semantic understanding through attention with efficient sequential processing via recurrence, reducing memory requirements during both training and inference. The model is trained from scratch on diverse open-source conversational and instruction-following datasets, followed by comprehensive evaluation against established baselines using standard language modeling metrics. The work includes full production infrastructure validation through cloud-based deployments, automated testing pipelines, and cost-efficiency analysis. Results demonstrate that hybrid architectures can achieve competitive performance with reduced computational footprint, providing a viable alternative for developing specialized language models in resource-constrained environments.
\end{abstract}
\vspace*{\fill}
\clearpage

\vspace*{\fill}
\begin{abstract}
    Este trabajo investiga el diseño e implementación de un Small Language Model (SLM) utilizando una arquitectura híbrida que integra mecanismos de atención Transformer \textit{multi-head attention} con Gated Recurrent Units (GRU). En lugar de depender exclusivamente de arquitecturas basadas en atención, el enfoque propuesto combina la comprensión semántica global mediante atención con el procesamiento secuencial eficiente mediante recurrencia, reduciendo los requisitos de memoria durante entrenamiento e inferencia. El modelo se entrena desde cero sobre datasets públicos diversos de diálogos e instrucciones, seguido de evaluación exhaustiva contra líneas base establecidas utilizando métricas estándar de modelado del lenguaje. El trabajo incluye validación completa de infraestructura en producción mediante despliegues en la nube, pipelines de testing automatizados, y análisis de eficiencia económica. Los resultados demuestran que arquitecturas híbridas pueden alcanzar rendimiento competitivo con menor costo computacional, proporcionando una alternativa viable para desarrollar modelos de lenguaje especializados en entornos con recursos limitados.
\end{abstract}
\vspace*{\fill}
\clearpage

\vspace*{\fill}
\begin{abstract}
    Aquest treball investiga el disseny i implementació d'un Small Language Model (SLM) utilitzant una arquitectura híbrida que integra mecanismes d'atenció Transformer \textit{multi-head attention} amb Gated Recurrent Units (GRU). En lloc de dependre exclusivament d'arquitectures basades en atenció, l'enfocament proposat combina la comprensió semàntica global mitjançant atenció amb el processament seqüencial eficient mitjançant recurrència, reduint els requisits de memòria durant entrenament i inferència. El model s'entrena des de zero sobre datasets públics diversos de diàlegs i instruccions, seguit d'avaluació exhaustiva contra línies base establertes utilitzant mètriques estàndard de modelatge del llenguatge. El treball inclou validació completa d'infraestructura en producció mitjançant desplegaments al núvol, pipelines de testing automatitzats, i anàlisi d'eficiència econòmica. Els resultats demostren que arquitectures híbrides poden assolir rendiment competitiu amb menor cost computacional, proporcionant una alternativa viable per desenvolupar models de llenguatge especialitzats en entorns amb recursos limitats.
\end{abstract}
\vspace*{\fill}
\clearpage
